\section{Execution-Graph Defect Theory}
\label{sec:theory}

Let $N_l$ denote normalization, and absorb it into the attention and feed-forward branches,
\begin{equation}
  a_l(x)=A_l(N_l(x)),\qquad g_l(x)=F_l(N_l(x)).
\end{equation}
The sequential and parallel realizations of block $l$ are
\begin{align}
  S_l(x)&=x+a_l(x)+g_l(x+a_l(x)),\\
  P_l(x)&=x+a_l(x)+g_l(x).
\end{align}
Their local execution-graph defect is
\begin{equation}
  d_l(x)=S_l(x)-P_l(x)=g_l(x+a_l(x))-g_l(x).
  \label{eq:local-defect}
\end{equation}
Thus the graph rewrite removes exactly the interaction between the attention update and the normalized nonlinear FFN.

\begin{theorem}[Exact interaction identity]
\label{thm:exact-defect}
If $g_l$ is continuously differentiable on $\{x+t a_l(x):t\in[0,1]\}$, then
\begin{equation}
  d_l(x)=\int_0^1 Jg_l(x+t a_l(x))a_l(x)\,dt.
  \label{eq:integral-defect}
\end{equation}
Consequently,
\begin{equation}
  \|d_l(x)\|\le
  \sup_{t\in[0,1]}\|Jg_l(x+t a_l(x))\|_{\mathrm{op}}\,\|a_l(x)\|.
\end{equation}
\end{theorem}

\begin{proof}
Fix $x$ and define $\gamma(t)=g_l(x+t a_l(x))$. The chain rule gives $\gamma'(t)=Jg_l(x+t a_l(x))a_l(x)$. Applying the fundamental theorem of calculus to $\gamma(1)-\gamma(0)$ proves Eq.~\eqref{eq:integral-defect}; the norm bound follows from the triangle inequality.
\end{proof}

\begin{theorem}[First-order defect approximation]
\label{thm:taylor-defect}
Suppose $Jg_l$ is $H_l$-Lipschitz on the segment in Theorem~\ref{thm:exact-defect}. Then
\begin{equation}
  d_l(x)=Jg_l(x)a_l(x)+r_l(x),
\end{equation}
where
\begin{equation}
  \|r_l(x)\|\le \frac{H_l}{2}\|a_l(x)\|^2.
\end{equation}
\end{theorem}

\begin{proof}
Subtract $Jg_l(x)a_l(x)$ from Eq.~\eqref{eq:integral-defect} and use Jacobian Lipschitzness:
\begin{align}
  \|r_l(x)\|
  &\le \int_0^1
  \|Jg_l(x+t a_l(x))-Jg_l(x)\|_{\mathrm{op}}\|a_l(x)\|\,dt\\
  &\le \int_0^1 H_l t\|a_l(x)\|^2\,dt
  =\frac{H_l}{2}\|a_l(x)\|^2.
\end{align}
\end{proof}

Let $m\in\{0,1\}^L$ denote an execution graph, where $m_l=1$ selects $P_l$ and $m_l=0$ selects $S_l$. Let $x_l^m$ and $x_l^0$ be the corresponding graph and all-sequential trajectories.

\begin{theorem}[Multi-layer graph-rewrite bound]
\label{thm:graph-bound}
Assume $S_l$ is $K_l$-Lipschitz on a domain containing both trajectories and $\sup_x\|d_l(x)\|\le D_l$ on that domain. Then
\begin{equation}
  \|x_L^m-x_L^0\|
  \le
  \sum_{l=0}^{L-1}m_lD_l\prod_{j=l+1}^{L-1}K_j.
  \label{eq:graph-bound}
\end{equation}
\end{theorem}

\begin{proof}
Let $e_l=\|x_l^m-x_l^0\|$. Adding and subtracting $S_l(x_l^m)$ gives
\begin{align}
  e_{l+1}
  &\le \|S_l(x_l^m)-S_l(x_l^0)\|
  +m_l\|P_l(x_l^m)-S_l(x_l^m)\|\\
  &\le K_l e_l+m_lD_l.
\end{align}
Unrolling from $e_0=0$ proves Eq.~\eqref{eq:graph-bound}.
\end{proof}

\begin{theorem}[Output-distribution bound]
\label{thm:kl-bound}
Let the final logit map $h$ be $K_h$-Lipschitz, and define $p_m=\mathrm{softmax}(h(x_L^m))$ and $p_0=\mathrm{softmax}(h(x_L^0))$. Then
\begin{equation}
  \mathrm{KL}(p_0\|p_m)
  \le \frac{K_h^2}{4}
  \left(
    \sum_{l=0}^{L-1}m_lD_l\prod_{j=l+1}^{L-1}K_j
  \right)^2.
  \label{eq:kl-bound}
\end{equation}
\end{theorem}

\begin{proof}
For $A(z)=\log\sum_i\exp z_i$, categorical KL is the Bregman divergence $A(z')-A(z)-\nabla A(z)^\top(z'-z)$. The Hessian $\nabla^2A(z)=\mathrm{diag}(p)-pp^\top$ has Euclidean operator norm at most $1/2$, so the divergence is at most $\frac14\|z'-z\|^2$. Apply the Lipschitz bound for $h$ followed by Theorem~\ref{thm:graph-bound}.
\end{proof}

\paragraph{First-order graph composition.}
Extend the binary mask to $u\in[0,1]^L$ using $B_l^u(x)=S_l(x)-u_ld_l(x)$. A multivariate Taylor expansion around the all-sequential graph gives
\begin{equation}
  z(u)-z(0)=\sum_lu_lv_l+R(u),
  \qquad \|R(u)\|\le C\|u\|_1^2,
\end{equation}
where $v_l=\partial z/\partial u_l|_{u=0}$ is the propagated single-layer defect. Quadratic output metrics contain pairwise terms $v_i^THv_j$; weak or negative average cross terms yield approximate additivity or subadditivity. Empirically, across both corpora we find
\begin{equation}
  D(m)\approx \alpha\sum_{l:m_l=1}D(e_l),
  \qquad \alpha\in[0.69,0.72].
\end{equation}
The coefficient is an empirical calibration parameter, not a universal constant.

\paragraph{Certificate boundary.}
Theorems~\ref{thm:exact-defect}--\ref{thm:kl-bound} require segment or trajectory-domain bounds. A defect measured only on a reference prompt is not automatically a uniform $D_l$. Our aggregate graph calibration transfers well across corpora, but prompt-adaptive selection exhibits tail undercoverage on held-out ClimbMix prompts. We therefore use \emph{defect predictor} or \emph{budgeted graph ranker} for the empirical method and reserve \emph{certificate} for settings with validated coverage.

\paragraph{Fisher decomposition of multi-layer KL.}
The first-order expansion $z(m)-z(0)\approx\sum_{l\in m}v_l$ and the quadratic structure of KL divergence near the baseline yield
\begin{equation}
  D(m)
  =\sum_{l\in m}d_l+\sum_{\substack{i<j\\i,j\in m}}C_{ij}+O(\|u\|_1^3),
  \label{eq:fisher-decompose}
\end{equation}
where $d_l=\tfrac12 v_l^\top\!Fv_l$ with $F=\mathrm{diag}(p_0)-p_0p_0^\top$ the Fisher information of the sequential output and $C_{ij}=v_i^\top\!Fv_j$.

\begin{proposition}[Balanced cancellation and the composition exponent]
\label{prop:balanced-cancel}
Let $\eta=D(m_{\mathrm{all}})/\!\sum_l d_l$ denote the all-parallel ratio, $L$ the number of layers, and $\bar C=\sum_{i<j}C_{ij}\big/\!\binom{L}{2}$ the mean pairwise Fisher cross-term.
\begin{enumerate}[label=(\alph*)]
\item Under the uniform cross-term approximation $C_{ij}\approx\bar C$ for all pairs, the composition ratio for a $k$-layer subset is
\begin{equation}
  r(k)\;=\;\frac{D(m)}{D_{\mathrm{add}}(m)}
  \;\approx\;1-\frac{(k-1)(1-\eta)}{L-1}.
  \label{eq:ratio-linear}
\end{equation}
\item The fitted power-law exponent $\beta$ in $D(m)\approx\alpha_0\!\sum_l d_l\cdot|m|^{-\beta}$ satisfies
\begin{equation}
  \beta\;\ge\;\beta_{\mathrm{BC}}
  \;=\;\frac{\log(1/\eta)}{\log L}.
  \label{eq:beta-lower}
\end{equation}
\end{enumerate}
\end{proposition}

\begin{proof}
(a)~The quadratic decomposition~\eqref{eq:fisher-decompose} with uniform $C_{ij}=\bar C$ gives $D(m)=\sum_{l\in m}d_l+\binom{k}{2}\bar C$. Setting $k=L$ yields $\bar C=(\eta-1)\sum_l d_l\big/\binom{L}{2}$. For a $k$-element subset with $\sum_{l\in m}d_l\approx(k/L)\sum_l d_l$, substitution and simplification give~\eqref{eq:ratio-linear}.

(b)~The linear model~\eqref{eq:ratio-linear} satisfies $r(1)=1$ and $r(L)=\eta$. The empirical ratio curves are concave (faster initial decline, then saturation), so the best-fit power-law exponent through the same endpoints is at least $-\!\log\eta/\log L$.
\end{proof}

\paragraph{Why subadditivity is not a geometric necessity.}
Eq.~\eqref{eq:ratio-linear} shows that the sign of the composition correction depends on whether $\eta<1$ (subadditive) or $\eta>1$ (superadditive). Polymorphic-trained models achieve $\eta\in[0.40,0.64]$ across scales, while the sequential-trained 430M baseline has $\eta\approx1.23$---mildly superadditive composition ($\beta<0$). Subadditivity thus arises from the consistency objective driving the all-parallel logits toward the sequential logits, not from defect-vector geometry.

\paragraph{Role of defect orthogonality.}
The near-zero Euclidean cosines $\cos(\delta_i,\delta_j)\approx0$ reflect the high ambient dimension $d\gg L$ rather than a special geometric arrangement. They imply near-quadratic combination in hidden-state norm ($\|\sum v_l\|^2\approx\sum\|v_l\|^2$), strengthening the Lipschitz bound of Theorem~\ref{thm:graph-bound}. However, the KL cross-term $C_{ij}=v_i^\top\!Fv_j$ involves the Fisher metric $F\ne cI$; its sign is unconstrained by the Euclidean cosine.

\paragraph{Sources of excess $\beta$.}
The lower bound $\beta_{\mathrm{BC}}$ underestimates the fitted $\beta$ by a factor of 1.3--1.8 across our eight configurations. Two mechanisms account for the gap:
\begin{itemize}
\item \emph{Nonlinear remainder.} Measured pairwise interactions are approximately $5\times$ larger than the uniform cross-term prediction, indicating substantial $O(\|u\|_1^3)$ contributions that steepen the power-law fit.
\item \emph{Defect heterogeneity.} The coefficient of variation $\mathrm{CV}(d_l)$ correlates with the excess $\beta_{\mathrm{fit}}-\beta_{\mathrm{BC}}$ at $\rho=0.90$ ($p=0.002$). Heterogeneous single-layer effects amplify apparent subadditivity because masks selecting high-defect layers have disproportionately large $D_{\mathrm{add}}$ but only moderately larger $D$.
\end{itemize}
